% ============================================================
% Margin Requirements — Formula Reference
% Trade Republic — Margin Trading Methodology
% Compile with: pdflatex margin_methodology.tex
% ============================================================

\documentclass[9pt, a4paper]{extarticle}

% ── Packages ─────────────────────────────────────────────
\usepackage[a4paper, top=2.5cm, bottom=2.5cm, left=3cm, right=3cm]{geometry}
\usepackage{amsmath, amssymb, amsthm}
\usepackage{booktabs}
\usepackage{array}
\usepackage{tabularx}
\usepackage{titlesec}
\usepackage{enumitem}
\usepackage{fancyhdr}
\usepackage{hyperref}
\usepackage{microtype}
\usepackage[T1]{fontenc}
\usepackage[sfdefault]{inter}
\usepackage{parskip}
\usepackage{longtable}
\usepackage{textcomp}

% ── Hyperref ─────────────────────────────────────────────
\hypersetup{
    hidelinks,
    pdftitle  = {Margin Requirements — Formula Reference},
    pdfauthor = {},
}

% ── Section heading styles ────────────────────────────────
\titleformat{\section}
    {\large\bfseries}
    {\thesection.}{0.6em}{}
    [\titlerule]

\titleformat{\subsection}
    {\normalsize\bfseries}
    {\thesubsection}{0.6em}{}

\titleformat{\subsubsection}
    {\normalsize\itshape}
    {\thesubsubsection}{0.6em}{}

\titlespacing{\section}{0pt}{1.6em}{0.8em}
\titlespacing{\subsection}{0pt}{1.2em}{0.4em}
\titlespacing{\subsubsection}{0pt}{1.0em}{0.3em}

% ── Header / Footer ──────────────────────────────────────
\pagestyle{fancy}
\fancyhf{}
\renewcommand{\headrulewidth}{0.4pt}
\fancyhead[L]{\small\textit{Margin Requirements — Formula Reference}}
\fancyhead[R]{\small\textsc{Strictly Private \& Confidential} $|$ March 2026}
\fancyfoot[C]{\small\thepage}

% ── Formula block environment ─────────────────────────────
\newenvironment{formulabox}{\par\addvspace{4pt}}{\par\addvspace{4pt}}

% ── Variable table command ────────────────────────────────
\newcommand{\vartableheader}{%
    \textbf{Symbol} & \textbf{Description}\\
}

% ── Euro symbol ──────────────────────────────────────────
\newcommand{\EUR}[1]{\ifmmode\text{\texteuro}#1\else\texteuro#1\fi}

% ── Misc ─────────────────────────────────────────────────
\setlength{\parindent}{0pt}
\setlength{\parskip}{5pt}

% ─────────────────────────────────────────────────────────
\begin{document}
% ─────────────────────────────────────────────────────────

% ── Title page ───────────────────────────────────────────
\begin{titlepage}
    \centering
    \vspace*{4cm}
    {\rule{\linewidth}{1.2pt}}\\[0.6em]
    {\Huge\bfseries Margin Requirements}\\[0.4em]
    {\Large Formula Reference Document}\\[0.4em]
    {\rule{\linewidth}{1.2pt}}\\[2em]
    {\large Trade Republic — Margin Trading Methodology}\\[0.5em]
    {\normalsize Version 1.0 \quad|\quad March 2026}\\[0.5em]
    {\normalsize Classification: Strictly Private \& Confidential}
    \vfill
    {\small This document sets out all quantitative formulas used in Trade Republic's
    Margin Requirements Methodology. Formulas are presented in the order in which
    they appear in the methodology document, grouped by component. Variable
    definitions accompany each formula block. A consolidated formula reference
    appears in the Appendix.}
\end{titlepage}

\tableofcontents
\newpage

% ─────────────────────────────────────────────────────────
%  SECTION D — MARGIN
% ─────────────────────────────────────────────────────────
\section{Margin (Section D)}
\label{sec:margin}

\subsection{D.1 — Four Layers of Margin}

The margin requirement framework consists of four layers applied on an
instrument-by-instrument basis:

\begin{enumerate}[leftmargin=2em]
    \item \textbf{Initial Margin (IM)} — minimum equity to open a leveraged position.
    \item \textbf{Maintenance Margin (MM)} — minimum equity the customer must maintain
          at all times after a position is opened.
    \item \textbf{Variation Margin (VM)} — daily cash P\&L settlement applied to all
          open positions subject to mark-to-market valuation.
    \item \textbf{Overnight Maintenance Margin} — additional haircut-based collateral
          requirement applied at end-of-day to account for the inability to monitor
          positions outside of trading hours.
\end{enumerate}

\subsubsection{D.1.1 — Initial Margin Rate}

The initial margin rate is the minimum equity the customer must post before a
leveraged position can be opened. It is calibrated at instrument level using a
scenario-based Value-at-Risk methodology and is expressed as a fraction of the
notional position value. The rate is composed of two additive components: a Market
Risk Component (MRC), which captures the potential loss over the assumed liquidation
horizon at a 99\% confidence level, and a Liquidity Component (LC), which accounts
for the cost of unwinding the position in adverse market conditions.

\begin{formulabox}
\textbf{Initial Margin Rate}
\begin{equation*}
    \mathrm{IM}_{\text{rate}}(i) \;=\; \mathrm{MRC}(i) \;+\; \mathrm{LC}(i)
\end{equation*}
\textbf{Market Risk Component}
\begin{equation*}
    \mathrm{MRC}(i) \;=\; \max\!\Bigl(\mathrm{RVaR}^{\text{hist}}(i) + \mathrm{CB}(i),\;
                           w_s \cdot \mathrm{RVaR}^{\text{stress}}(i)\Bigr)
\end{equation*}
\textbf{Liquidity Component}
\begin{equation*}
    \mathrm{LC}(i) \;=\; \mathrm{DF} \cdot \mathrm{RVaR}^{\text{hist}}(i) \cdot \lambda(i)
\end{equation*}
\end{formulabox}

\begin{tabularx}{\linewidth}{lX}
\toprule
\vartableheader
\midrule
$\mathrm{IM}_{\text{rate}}(i)$      & Initial margin rate for instrument $i$ as a fraction of position value \\
$\mathrm{MRC}(i)$                   & Market Risk Component for instrument $i$ \\
$\mathrm{LC}(i)$                    & Liquidity Component for instrument $i$ \\
$\mathrm{RVaR}^{\text{hist}}(i)$    & Robust VaR from filtered historical scenarios (see Section~D.2.4) \\
$\mathrm{RVaR}^{\text{stress}}(i)$  & Robust VaR from stressed scenarios (see Section~D.2.2) \\
$\mathrm{CB}(i)$                    & Correlation Break adjustment (see Section~D.2.4) \\
$w_s$                               & Stress weight (default: 0.5) \\
$\mathrm{DF}$                       & Diversification factor for the liquidation group \\
$\lambda(i)$                        & Days-to-liquidate scaling factor for instrument $i$ \\
\bottomrule
\end{tabularx}

The initial margin rate is computed on a standalone basis per instrument, meaning
each instrument is treated as its own liquidation group. This ensures that the rate
displayed to the customer prior to trade execution reflects the true cost of holding
that specific instrument on leverage, without benefit from diversification across
instruments the customer does not hold. Full details of the calibration methodology
are provided in Section~D.2.

\subsubsection{D.1.2 — Maintenance Margin Rate}

The maintenance margin rate represents the minimum equity ratio the customer must
maintain at all times after a leveraged position is opened. It is expressed as a
minimum Adjusted Collateral Value relative to the outstanding loan, governed by the
LTV monitoring threshold $\overline{\mathrm{LTV}}$. For each instrument $i$, the
required minimum ACV at time $t$ is:

\begin{formulabox}
\textbf{Minimum Required ACV}
\begin{equation*}
    \mathrm{ACV}^{\min}_{t} \;=\; \frac{L_t}{\overline{\mathrm{LTV}}}
\end{equation*}
\textbf{Maintenance Margin Rate}
\begin{equation*}
    \mathrm{MM}_{\text{rate}}(i) \;=\; 1 \;-\; \overline{\mathrm{LTV}} \cdot \bigl(1 - h_{\mathrm{eff}}(i)\bigr)
\end{equation*}
\end{formulabox}

\begin{tabularx}{\linewidth}{lX}
\toprule
\vartableheader
\midrule
$\mathrm{ACV}^{\min}_{t}$   & Minimum Adjusted Collateral Value required at time $t$ \\
$L_t$                        & Outstanding loan balance at time $t$ \\
$\overline{\mathrm{LTV}}$    & LTV monitoring threshold (default: 75\%) \\
$h_{\mathrm{eff}}(i)$        & Effective haircut for instrument $i$ (see Section~E.2) \\
$\mathrm{MM}_{\text{rate}}(i)$ & Maintenance margin rate for instrument $i$ as a fraction of market value \\
\bottomrule
\end{tabularx}

The maintenance margin rate is instrument-specific: instruments with higher effective
haircuts require a larger minimum equity cushion for the same loan amount. A customer
whose ACV falls below $\mathrm{ACV}^{\min}_t$ — equivalently, whose LTV exceeds
$\overline{\mathrm{LTV}}$ — is subject to the escalation framework described in
Section~D.1.4.

\paragraph{Two-Phase LTV Framework.}
The margin framework operates in two distinct phases tied to the liquidation horizon
$T_{\text{liq}}$ (default: 4 trading days):

\begin{itemize}[leftmargin=2em]
    \item \textbf{Phase~1 — Opening ($t = 0$):} At trade inception the customer must
          post Initial Margin, so the opening LTV is directly determined by the IM
          rate:
          \begin{equation*}
              \mathrm{LTV}_0 \;=\; \frac{L_0}{\mathrm{ACV}_0}
                              \;=\; \frac{1 - \mathrm{IM}_{\text{rate}}(i)}{1 - h_{\mathrm{eff}}(i)}
          \end{equation*}

    \item \textbf{Phase~2 — Seasoned ($t > T_{\text{liq}}$):} Once the position has
          been held beyond the liquidation horizon, the Initial Margin requirement no
          longer applies. The binding constraint becomes the maintenance margin
          threshold $\overline{\mathrm{LTV}} = 75\%$, monitored continuously via
          daily Variation Margin and end-of-day Overnight Maintenance Margin. The
          outstanding loan evolves as:
          \begin{equation*}
              L_t \;=\; L_0 - \sum_{\tau=1}^{t} \mathrm{VM}_\tau
          \end{equation*}
          and a margin event is triggered whenever $\mathrm{LTV}_t > \overline{\mathrm{LTV}}$.
\end{itemize}

This two-phase structure ensures that the Initial Margin covers the worst-case loss
over the full liquidation horizon, while the lighter Maintenance Margin framework
governs the ongoing life of the position once that horizon has elapsed.

\subsubsection{D.1.3 — Variation Margin}

The profit or loss arising from the change in the theoretical value of an open
position between two consecutive trading days. Gains and losses are settled
symmetrically and variation margin is not held as collateral.

\begin{formulabox}
\textbf{Variation Margin}
\begin{equation*}
    \mathrm{VM}_{i,t} =
        \bigl(P_{i,t} - P_{i,t-1}\bigr) \times N_i
\end{equation*}
\textit{Positive VM is credited to the customer; negative VM is debited.
Applied to all positions on every trading day.}
\end{formulabox}

\begin{tabularx}{\linewidth}{lX}
\toprule
\vartableheader
\midrule
$\mathrm{VM}_{i,t}$  & Variation Margin for instrument $i$ on day $t$ \\
$P_{i,t}$            & Settlement price of instrument $i$ on day $t$ \\
$P_{i,t-1}$          & Settlement price of instrument $i$ on the previous trading day \\
$N_i$                & Quantity held of instrument $i$ \\
\bottomrule
\end{tabularx}

\subsubsection{D.1.4 — Overnight Maintenance Margin}

At the close of each trading day, Trade Republic applies an overnight maintenance
margin requirement to all open leveraged positions, reflecting the increased risk
arising from the inability to monitor or act on adverse price movements outside of
trading hours. The overnight margin is expressed as a minimum Adjusted Collateral
Value relative to the outstanding loan, calibrated using a one-day parametric
Value-at-Risk at the 99\% confidence level, scaled to the overnight holding period.
Specifically, the overnight haircut applied to each instrument is derived as:

\begin{formulabox}
\textbf{Overnight Haircut}
\begin{equation*}
    h_{\text{overnight}}(i) \;=\;
        h_{\text{floor}}\!\bigl(\text{class}(i)\bigr)
        \;+\;
        w_{\text{buf}} \cdot
        \max\!\Bigl(0,\; h_{\text{vol}}(i) - h_{\text{floor}}\!\bigl(\text{class}(i)\bigr)\Bigr)
\end{equation*}
\end{formulabox}

where $h_{\text{floor}}$ is the asset-class minimum haircut, $h_{\text{vol}}$ is the
volatility-implied haircut based on the instrument's EWMA historical volatility, and
$w_{\text{buf}}$ is a volatility buffer weight. This ensures that the collateral
discount applied overnight is anchored to a regulatory-grade asset class floor while
remaining responsive to current market conditions. Any customer whose LTV breaches
the monitoring threshold of 75\% at end-of-day is subject to an overnight margin
review, with corrective action --- either additional collateral posting or partial
position reduction --- required prior to the opening of the next trading session.

\begin{tabularx}{\linewidth}{lX}
\toprule
\vartableheader
\midrule
$h_{\text{overnight}}(i)$ & Overnight haircut for instrument $i$ \\
$h_{\text{floor}}(\text{class}(i))$ & Asset-class minimum haircut floor \\
$h_{\text{vol}}(i)$ & Volatility-implied haircut based on EWMA historical volatility \\
$w_{\text{buf}}$ & Volatility buffer weight $(0 \leq w_{\text{buf}} \leq 1)$ \\
\bottomrule
\end{tabularx}

\subsubsection{D.1.5 — Monitoring Framework}

The monitoring framework operates in two distinct phases determined by the
liquidation horizon $T_{\text{liq}}$ (default: 4 trading days). The set of active
checks and the associated breach actions differ between the two phases.

\paragraph{Phase 1 — Within Liquidation Horizon ($0 < t \leq T_{\text{liq}}$).}
During this phase the Initial Margin buffer is still the primary safeguard. Daily
Variation Margin is applied and the LTV is monitored continuously. A breach signals
that adverse price moves have consumed the IM buffer faster than the position can be
wound down.

\begin{center}
\begin{tabularx}{0.96\linewidth}{lXl}
\toprule
\textbf{Check} & \textbf{Condition} & \textbf{Action} \\
\midrule
IM\,+\,VM intraday
    & $\mathrm{LTV}_t = L_t / \mathrm{ACV}_t > \overline{\mathrm{LTV}}$
    & Monitoring alert \\
IM\,+\,VM end-of-day
    & $\mathrm{LTV}_t^{\text{overnight}} > \overline{\mathrm{LTV}}$
    & Overnight margin review \\
Liquidation
    & $\mathrm{LTV}_t > \mathrm{LTV}_{\text{liq}}$
    & Auto-liquidate \\
\bottomrule
\end{tabularx}
\end{center}

where the loan evolves daily as:
\[
    L_t = L_0 - \sum_{\tau=1}^{t} \mathrm{VM}_\tau
\]
and $\mathrm{LTV}_t^{\text{overnight}}$ uses the overnight haircut $h_{\text{overnight}}(i)$
(Section~D.1.4) in place of $h_{\text{eff}}(i)$ when computing $\mathrm{ACV}_t$.

\paragraph{Phase 2 — Seasoned Position ($t > T_{\text{liq}}$).}
Once the position has been held beyond the liquidation horizon, the Initial Margin
requirement is no longer binding. The lighter Maintenance Margin framework applies,
supplemented by the overnight check at each close of business.

\begin{center}
\begin{tabularx}{0.96\linewidth}{lXl}
\toprule
\textbf{Check} & \textbf{Condition} & \textbf{Action} \\
\midrule
VM intraday
    & $\mathrm{LTV}_t > \overline{\mathrm{LTV}}$
    & Monitoring alert \\
MM end-of-day
    & $\mathrm{LTV}_t^{\text{overnight}} > \overline{\mathrm{LTV}}$
    & Overnight margin review \\
Liquidation
    & $\mathrm{LTV}_t > \mathrm{LTV}_{\text{liq}}$
    & Auto-liquidate \\
\bottomrule
\end{tabularx}
\end{center}

\paragraph{Escalation Logic.}
The three breach conditions are ordered by severity and share the same LTV
thresholds across both phases:

\begin{center}
\begin{tabularx}{0.96\linewidth}{lllX}
\toprule
\textbf{Level} & \textbf{Trigger} & \textbf{LTV Threshold} & \textbf{Description} \\
\midrule
1 — Monitoring
    & $\mathrm{ACV}_t < L_t\,/\,\overline{\mathrm{LTV}}$
    & $75\%$
    & Enhanced surveillance; no immediate customer action required \\
2 — Overnight breach
    & $\mathrm{ACV}_t^{\text{overnight}} < L_t\,/\,\overline{\mathrm{LTV}}$
    & $75\%$ (overnight ACV)
    & Customer must post collateral or reduce loan before next open \\
3 — Liquidation
    & $\mathrm{ACV}_t < L_t\,/\,\mathrm{LTV}_{\text{liq}}$
    & $90\%$
    & Positions auto-liquidated to repay outstanding loan \\
\bottomrule
\end{tabularx}
\end{center}

\begin{tabularx}{\linewidth}{lX}
\toprule
\vartableheader
\midrule
$T_{\text{liq}}$                    & Liquidation horizon (default: 4 trading days) \\
$L_t$                               & Outstanding loan at time $t$, net of accumulated VM \\
$\mathrm{ACV}_t$                    & Adjusted Collateral Value at time $t$ using intraday haircut $h_{\text{eff}}$ \\
$\mathrm{ACV}_t^{\text{overnight}}$ & Adjusted Collateral Value at close using overnight haircut $h_{\text{overnight}}$ \\
$\overline{\mathrm{LTV}}$           & Monitoring threshold (default: $75\%$) \\
$\mathrm{LTV}_{\text{liq}}$         & Liquidation threshold (default: $90\%$) \\
\bottomrule
\end{tabularx}

\newpage

% ─────────────────────────────────────────────────────────
\subsection{D.2 — Calibration of Margin Rates}

Margin rates are calibrated at instrument level using a Value-at-Risk (VaR)
methodology designed to estimate the loss that could arise during the assumed
liquidation horizon at the chosen confidence level.

\subsubsection{D.2.1 — Mapping Liquidation Groups}

Every instrument is assigned to a liquidation group $l$ whose members share
sufficient characteristics — market, asset class, and risk profile — to be
auctioned together. Each liquidation group may be further divided into
liquidation group \textit{splits}.

% ─────────────────────────────────────────────────────────
\subsubsection{D.2.2 — Scenario Construction}

Two distinct scenario sets are constructed to generate hypothetical future
prices for every instrument.

\paragraph{Filtered Historical Scenarios.}
The first set consists of 750 observations drawn from recent market history.
For each historical date $s$, the realised return over the relevant holding
period is rescaled to reflect current volatility conditions — ensuring that
calm historical periods do not produce an underestimate of today's risk:

\begin{formulabox}
\textbf{Filtered (Volatility-Scaled) Return}
\begin{equation*}
    r_{\mathrm{filt}}(i,s)
        = r_{\mathrm{hist}}(i,s)
          \times \frac{\hat{\sigma}_{i,0}}{\hat{\sigma}_{i,s}}
\end{equation*}
\end{formulabox}

\begin{formulabox}
\textbf{Filtered Historical Scenario Price}
\begin{equation*}
    P_{\mathrm{filt}}(i,s)
        = P_{i,0} \times \bigl(1 + r_{\mathrm{filt}}(i,s)\bigr)
\end{equation*}
\end{formulabox}

\begin{tabularx}{\linewidth}{lX}
\toprule
\vartableheader
\midrule
$r_{\mathrm{hist}}(i,s)$   & Historical return of instrument $i$ over the holding period ending on date $s$ \\
$\hat{\sigma}_{i,0}$        & Current volatility estimate for instrument $i$ (at $T=0$) \\
$\hat{\sigma}_{i,s}$        & Historical volatility estimate for instrument $i$ at the time of observation $s$ \\
$P_{i,0}$                   & Current market price of instrument $i$ (at $T=0$) \\
$P_{\mathrm{filt}}(i,s)$   & Filtered scenario price for instrument $i$ under scenario $s$ \\
\bottomrule
\end{tabularx}

\paragraph{Stressed Scenarios.}
The second set consists of 250 scenarios drawn from identified stress periods in
market history (exceptional volatility, sharp correlation dislocations, extreme
price movements). Unlike the filtered scenarios, stressed returns are applied
\emph{without} any volatility adjustment:

\begin{formulabox}
\textbf{Stressed Scenario Price}
\begin{equation*}
    P_{\mathrm{stress}}(i,s)
        = P_{i,0} \times \bigl(1 + r_{\mathrm{stress}}(i,s)\bigr)
\end{equation*}
\end{formulabox}

\begin{tabularx}{\linewidth}{lX}
\toprule
\vartableheader
\midrule
$r_{\mathrm{stress}}(i,s)$  & Raw historical return of instrument $i$ from a stress period, scenario $s$ \\
$P_{\mathrm{stress}}(i,s)$  & Stressed scenario price for instrument $i$ under scenario $s$ \\
\bottomrule
\end{tabularx}

% ─────────────────────────────────────────────────────────
\subsubsection{D.2.3 — Profit and Loss Distributions}

For each scenario, position-level P\&L is aggregated across instruments within
a liquidation group split to produce the P\&L distribution used for VaR estimation.

\begin{formulabox}
\textbf{Instrument P\&L under Scenario $s$}
\begin{equation*}
    \mathrm{PL}(i,s) =
        \bigl(P_{i,0} - P_{\mathrm{scenario}}(i,s)\bigr) \times N_i
\end{equation*}
\textit{$P_{\mathrm{scenario}}$ is either $P_{\mathrm{filt}}$ or $P_{\mathrm{stress}}$
depending on the scenario sub-sample.}
\end{formulabox}

\begin{formulabox}
\textbf{Liquidation Group Split P\&L}
\begin{equation*}
    \mathrm{PL}(l,s) = \sum_{i \,\in\, l} \mathrm{PL}(i,s)
\end{equation*}
\end{formulabox}

\begin{tabularx}{\linewidth}{lX}
\toprule
\vartableheader
\midrule
$\mathrm{PL}(i,s)$  & P\&L of instrument $i$ under scenario $s$ \\
$\mathrm{PL}(l,s)$  & Portfolio P\&L for liquidation group split $l$ under scenario $s$ \\
$N_i$               & Quantity held of instrument $i$ \\
$i \in l$           & All instruments assigned to liquidation group split $l$ \\
\bottomrule
\end{tabularx}

\newpage

% ─────────────────────────────────────────────────────────
\subsubsection{D.2.4 — Market Risk Component}

\paragraph{D.2.4.1 — Market Risk Model.}

The market risk component is derived from VaR figures computed from the P\&L
distributions at a target confidence level of $99\%$.

\medskip
\textbf{Empirical Quantile.}
Where the target quantile falls between two observations, linear interpolation
is applied. Let the P\&L vector be sorted in ascending order with $n$
observations; the quantile of the $i$-th sorted observation is:

\begin{formulabox}
\textbf{Empirical Quantile}
\begin{equation*}
    q(i) = 100\% \times \frac{i - 0.5}{n}
\end{equation*}
\end{formulabox}

\begin{formulabox}
\textbf{Interpolated VaR at confidence level $\alpha$}
\begin{equation*}
    \mathrm{VaR}(\alpha) = -\left[
        \frac{\mathrm{PL}(q^{+})\,\bigl(q^{+} - (1-\alpha)\cdot100\bigr)
            + \mathrm{PL}(q^{-})\,\bigl((1-\alpha)\cdot100 - q^{-}\bigr)}
             {q^{+} - q^{-}}
    \right]
\end{equation*}
\textit{$q^{+}$ and $q^{-}$ are the nearest empirical quantiles above and below
$(1-\alpha)\times100$ respectively. VaR is bounded from below at zero.}
\end{formulabox}

\medskip
\textbf{Robust VaR.}
To improve statistical robustness, VaR is first estimated at a lower
confidence level ($95\%$ for filtered historical; $90\%$ for stressed scenarios)
and then scaled to the $99\%$ target using quantiles of a heavy-tailed
Student-$t$ distribution:

\begin{formulabox}
\textbf{Robust VaR Scaling Factor}
\begin{equation*}
    \mathrm{SF} =
        \frac{t_{\nu}^{-1}(\alpha_{\mathrm{high}})}
             {t_{\nu}^{-1}(\alpha_{\mathrm{low}})}
\end{equation*}
\end{formulabox}

\begin{formulabox}
\textbf{Robust VaR}
\begin{equation*}
    \mathrm{RVaR}(l,a) =
        \mathrm{VaR}(\alpha_{\mathrm{low}},\,l,\,a) \times \mathrm{SF}
\end{equation*}
\end{formulabox}

\begin{tabularx}{\linewidth}{lX}
\toprule
\vartableheader
\midrule
$\alpha$                   & Target confidence level ($0.99$) \\
$\alpha_{\mathrm{high}}$   & Target confidence level ($0.99$) \\
$\alpha_{\mathrm{low}}$    & Base confidence level
                             ($0.95$ for filtered historical; $0.90$ for stressed) \\
$t_{\nu}^{-1}(\alpha)$     & Quantile (inverse CDF) of the Student-$t$ distribution
                             at confidence $\alpha$ with $\nu$ degrees of freedom \\
$\nu$                      & Degrees of freedom of the Student-$t$ distribution \\
$\mathrm{RVaR}(l,a)$      & Robust VaR for split $l$, scenario sub-sample $a$ \\
$n$                        & Number of observations in the scenario sub-sample \\
$q^{+},\,q^{-}$           & Nearest empirical quantiles bracketing the target percentile \\
\bottomrule
\end{tabularx}

\medskip
\textbf{Market Risk Component Aggregation.}
Robust VaR is computed separately for each scenario sub-sample. The mean across
sub-samples is taken within each scenario type. The final market risk component
is the maximum of the filtered historical mean (plus the Correlation Break
adjustment) and the stressed mean scaled by the stress weighting factor
$w_s \in [0.60,\,0.70]$:

\begin{formulabox}
\textbf{Market Risk Component}
\begin{equation*}
    \mathrm{MRC}(l) = \max\!\Biggl(
        \underbrace{\overline{\mathrm{RVaR}}_{\mathrm{hist}}(l)
                    + \overline{\mathrm{CB}}(l)}_{\text{filtered historical}},\;
        \underbrace{w_s \cdot
            \overline{\mathrm{RVaR}}_{\mathrm{stress}}(l)}_{\text{stressed}}
    \Biggr)
\end{equation*}
where $\overline{(\cdot)}$ denotes the mean across scenario sub-samples $a$.
\end{formulabox}

\begin{tabularx}{\linewidth}{lX}
\toprule
\vartableheader
\midrule
$\mathrm{MRC}(l)$                                    & Market Risk Component for split $l$ \\
$\overline{\mathrm{RVaR}}_{\mathrm{hist}}(l)$        & Mean Robust VaR — filtered historical scenarios \\
$\overline{\mathrm{RVaR}}_{\mathrm{stress}}(l)$      & Mean Robust VaR — stressed scenarios \\
$\overline{\mathrm{CB}}(l)$                          & Mean Correlation Break Adjustment \\
$w_s$                                                & Stressed scenario weighting factor (typically $0.60$–$0.70$) \\
\bottomrule
\end{tabularx}

\newpage

% ─────────────────────────────────────────────────────────
\paragraph{D.2.4.2 — Model Risk Buffer (Correlation Break Adjustment).}

A Correlation Break adjustment is applied to the market risk component to
account for the risk that correlations shift adversely during a liquidation.
The VaR is recalculated over rolling time-windows within the historical data;
excess risk relative to the full-sample estimate is aggregated via a
root-mean-square measure and clipped between a floor $\Phi$ and cap $\psi$
(both expressed as fractions of $\mathrm{RVaR}_{\mathrm{hist}}$):

\begin{formulabox}
\textbf{Excess Risk per Rolling Sub-Window $T_j$}
\begin{equation*}
    \mathrm{RVaR}^{+}(l,a,T_j) =
        \max\!\Bigl(0,\;
            \mathrm{RVaR}_{\mathrm{sub}}(l,a,T_j)
            - \mathrm{RVaR}_{\mathrm{hist}}(l,a)
        \Bigr)
\end{equation*}
\end{formulabox}

\begin{formulabox}
\textbf{Root Mean Square Excess Risk}
\begin{equation*}
    \mathrm{RMSE}(l,a) =
        \sqrt{\frac{1}{n^{+}_a}
              \sum_{j=1}^{n^{+}_a}
              \bigl[\mathrm{RVaR}^{+}(l,a,T_j)\bigr]^{2}}
\end{equation*}
\end{formulabox}

\begin{formulabox}
\textbf{Correlation Break Adjustment}
\begin{equation*}
    \mathrm{CB}(l,a) = \mathrm{clip}\!\Bigl(
        \kappa \cdot \mathrm{RMSE}(l,a),\;
        \Phi \cdot \mathrm{RVaR}_{\mathrm{hist}}(l,a),\;
        \psi \cdot \mathrm{RVaR}_{\mathrm{hist}}(l,a)
    \Bigr)
\end{equation*}
\end{formulabox}

\begin{tabularx}{\linewidth}{lX}
\toprule
\vartableheader
\midrule
$\mathrm{RVaR}_{\mathrm{sub}}(l,a,T_j)$ & Robust VaR over rolling window ending at $T_j$,
                                           sub-sample $a$ \\
$\mathrm{RVaR}_{\mathrm{hist}}(l,a)$    & Full sub-sample Robust VaR for split $l$,
                                           sub-sample $a$ \\
$\mathrm{RVaR}^{+}(l,a,T_j)$           & Excess risk for rolling window ending at $T_j$ \\
$\mathrm{RMSE}(l,a)$                    & Root mean square of positive excess risk figures \\
$\kappa$                                & Calibration scaling factor \\
$n^{+}_a$                               & Count of sub-windows with positive excess risk \\
$\Phi$                                  & Floor as fraction of $\mathrm{RVaR}_{\mathrm{hist}}$
                                           (minimum charge) \\
$\psi$                                  & Cap as fraction of $\mathrm{RVaR}_{\mathrm{hist}}$
                                           (maximum charge) \\
$\mathrm{CB}(l,a)$                      & Correlation Break Adjustment for split $l$,
                                           sub-sample $a$ \\
\bottomrule
\end{tabularx}

\newpage

% ─────────────────────────────────────────────────────────
\subsubsection{D.2.5 — Liquidity Risk Component}

The liquidity risk component reflects the market impact and execution costs of
liquidating positions in real market conditions. It is computed at the position
level and then aggregated at the liquidation group split level with a
diversification benefit applied.

\begin{formulabox}
\textbf{Position-Level Liquidity Risk Component}
\begin{equation*}
    \mathrm{LC}(i) =
        \mathrm{RVaR}_{\mathrm{pos}}(i) \times \lambda(i)
\end{equation*}
\textit{The component scales with the positional market risk, reflecting that
volatile instruments carry greater execution risk. $\lambda(i)$ is the
instrument-level liquidity factor, capturing market impact relative to traded
volume.}
\end{formulabox}

\begin{formulabox}
\textbf{Diversification Factor}
\begin{equation*}
    \mathrm{DF}(l) =
        \max\!\left(
            \mathrm{DF}_{\mathrm{floor}},\;
            \frac{\mathrm{RVaR}_{\mathrm{portfolio}}(l)}
                 {\displaystyle\sum_{i \,\in\, l} \mathrm{RVaR}_{\mathrm{pos}}(i)}
        \right)
\end{equation*}
\textit{A well-diversified portfolio, where individual risks partially offset
each other, receives a proportional reduction. The floor ensures a baseline
liquidity charge is always present.}
\end{formulabox}

\begin{formulabox}
\textbf{Aggregate Liquidity Risk Component}
\begin{equation*}
    \mathrm{LC}(l) =
        \mathrm{DF}(l) \times \sum_{i \,\in\, l} \mathrm{LC}(i)
\end{equation*}
\end{formulabox}

\begin{tabularx}{\linewidth}{lX}
\toprule
\vartableheader
\midrule
$\mathrm{LC}(i)$                             & Position-level Liquidity Risk Component \\
$\mathrm{RVaR}_{\mathrm{pos}}(i)$           & Positional Robust VaR for instrument $i$
                                               $= \mathrm{RVaR}_{\mathrm{unit}}(i) \times |N_i|$ \\
$\lambda(i)$                                 & Liquidity factor for instrument $i$ \\
$\mathrm{DF}(l)$                             & Diversification Factor for split $l$ \\
$\mathrm{DF}_{\mathrm{floor}}$              & Minimum diversification factor \\
$\mathrm{RVaR}_{\mathrm{portfolio}}(l)$     & Portfolio-level Robust VaR for split $l$ \\
$\mathrm{LC}(l)$                             & Aggregate Liquidity Risk Component for split $l$ \\
\bottomrule
\end{tabularx}

% ─────────────────────────────────────────────────────────
\subsubsection{D.2.6 — Aggregation and Final Initial Margin}

\begin{formulabox}
\textbf{Initial Margin — Liquidation Group Split}
\begin{equation*}
    \mathrm{IM}(l) = \mathrm{MRC}(l) + \mathrm{LC}(l)
\end{equation*}
\end{formulabox}

\begin{formulabox}
\textbf{Total Initial Margin — Customer Portfolio}
\begin{equation*}
    \mathrm{IM}_{\mathrm{total}} = \sum_{l} \mathrm{IM}(l)
\end{equation*}
\textit{No offsets are recognised across liquidation group splits.}
\end{formulabox}

\begin{tabularx}{\linewidth}{lX}
\toprule
\vartableheader
\midrule
$\mathrm{IM}(l)$               & Initial Margin for liquidation group split $l$ \\
$\mathrm{MRC}(l)$              & Market Risk Component for split $l$ \\
$\mathrm{LC}(l)$               & Liquidity Risk Component for split $l$ \\
$\mathrm{IM}_{\mathrm{total}}$ & Total portfolio Initial Margin \\
\bottomrule
\end{tabularx}

\newpage

% ─────────────────────────────────────────────────────────
%  SECTION E — LTV CALCULATION METHODOLOGY
% ─────────────────────────────────────────────────────────
\section{LTV Calculation Methodology (Section E)}
\label{sec:ltv}

\subsection{E.2 — Collateral LTV: Haircuts on Posted Collateral}

\subsubsection{E.2.2 — Haircut Calibration}

For each eligible collateral instrument, the haircut $H(i)$ is calibrated to
cover the expected price decline over the collateral liquidation horizon at a
$99\%$ confidence level, using a historical simulation consistent with the
initial margin scenario methodology. The resulting accepted value is the
fraction of market value recognised as satisfying the margin obligation.

\begin{formulabox}
\textbf{Collateral LTV}
\begin{equation*}
    \mathrm{LTV}_{\mathrm{coll}}(i) = 1 - H_{\mathrm{eff}}(i) 
\end{equation*}
\end{formulabox}

\begin{formulabox}
\textbf{Accepted Collateral Value}
\begin{equation*}
    \mathrm{ACV}(i) = \mathrm{MV}(i) \times \bigl(1 - H_{\mathrm{eff}}(i) \bigr)
\end{equation*}
\end{formulabox}

\begin{formulabox}
\textbf{Total Accepted Collateral Value}
\begin{equation*}
    \mathrm{ACV}_{\mathrm{total}} =
        \sum_{i} \mathrm{ACV}(i) + \sum_{j} \mathrm{Cash}(j)
\end{equation*}
\textit{Cash posted in EUR is accepted at full value ($H = 0$).
Non-EUR collateral carries an additional FX add-on applied to $H(i)$.}
\end{formulabox}

\begin{tabularx}{\linewidth}{lX}
\toprule
\vartableheader
\midrule
$H_{\mathrm{eff}}(i)$            & Haircut for collateral instrument $i$ \\
$\mathrm{MV}(i)$                & Current market value of collateral instrument $i$ \\
$\mathrm{LTV}_{\mathrm{coll}}(i)$ & Collateral LTV (accepted fraction of market value) \\
$\mathrm{ACV}(i)$               & Accepted Collateral Value of instrument $i$ \\
$\mathrm{ACV}_{\mathrm{total}}$ & Total Accepted Collateral Value \\
$\mathrm{Cash}(j)$              & Cash posted in currency $j$ \\
\bottomrule
\end{tabularx}

% ─────────────────────────────────────────────────────────
\subsubsection{E.2.3 — Concentration Add-ons}

Where a customer's collateral portfolio is concentrated within a correlated group
of instruments, an additional concentration haircut is applied. The add-on is
determined from the average pairwise correlation among instruments in each
collateral asset group $g$:

\begin{formulabox}
\textbf{Average Pairwise Correlation}
\begin{equation*}
    \bar{\rho}(g) =
        \frac{2}{m(m-1)}
        \sum_{i < j,\; i,j \,\in\, g} \rho(i,j)
\end{equation*}
\textit{Computed over all pairs $(i,j)$ within collateral asset group $g$;
$m$ is the number of instruments in the group.}
\end{formulabox}

\begin{formulabox}
\textbf{Diversification Score}
\begin{equation*}
    \mathrm{DS}(g) =
        \sqrt{\frac{1 + (m-1)\,\bar{\rho}(g)}{m}}
\end{equation*}
\textit{$\mathrm{DS}(g) \to 1/\sqrt{m}$ for uncorrelated instruments;
$\mathrm{DS}(g) \to 1$ as $\bar{\rho}(g) \to 1$ (fully correlated).
A fully uncorrelated portfolio receives no concentration add-on.}
\end{formulabox}

\begin{formulabox}
\textbf{Concentration Add-on Haircut}
\begin{equation*}
    H_{\mathrm{conc}}(i,g) =
        H_{\mathrm{base}}(i) \times
        \frac{\mathrm{DS}(g) - \mathrm{DS}_{\mathrm{floor}}}
             {1 - \mathrm{DS}_{\mathrm{floor}}}
\end{equation*}
\textit{Applied as an increment to the base haircut where
$\mathrm{DS}(g) > \mathrm{DS}_{\mathrm{floor}}$.
As average pairwise correlations rise, the diversification score falls and
the add-on increases proportionally.}
\end{formulabox}

\begin{formulabox}
\textbf{Effective Haircut}
\begin{equation*}
    H_{\mathrm{eff}}(i) = H_{\mathrm{base}}(i) + H_{\mathrm{conc}}(i,g)
\end{equation*}
\end{formulabox}

\begin{tabularx}{\linewidth}{lX}
\toprule
\vartableheader
\midrule
$\rho(i,j)$                      & Pairwise return correlation between instruments $i$ and $j$ \\
$\bar{\rho}(g)$                  & Average pairwise correlation within collateral asset group $g$ \\
$m$                              & Number of instruments in asset group $g$ \\
$\mathrm{DS}(g)$                 & Diversification Score for asset group $g$ \\
$\mathrm{DS}_{\mathrm{floor}}$   & Minimum diversification score (floor below which no add-on) \\
$H_{\mathrm{base}}(i)$           & Base haircut for collateral instrument $i$ \\
$H_{\mathrm{conc}}(i,g)$         & Concentration add-on haircut for instrument $i$ in group $g$ \\
$H_{\mathrm{eff}}(i)$            & Effective total haircut for instrument $i$ \\
\bottomrule
\end{tabularx}

\newpage

% ─────────────────────────────────────────────────────────
\subsection{E.3 — Loan-to-Value for Margin Lending and Leveraged Positions}

\subsubsection{E.3.2 — LTV Calculation}

The LTV ratio measures the relationship between the outstanding loan exposure and
the haircut-adjusted collateral value at any point in time. LTV ratios are
calculated in real time using current market prices, with foreign currency
positions converted to EUR at current spot rates.

\begin{formulabox}
\textbf{Net Liquidation Value}
\begin{equation*}
    \mathrm{NLV}_t = \mathrm{MV}_{\mathrm{portfolio},t} - L_t
\end{equation*}
\end{formulabox}

\begin{formulabox}
\textbf{LTV Ratio}
\begin{equation*}
    \mathrm{LTV}_t =
        \frac{L_t}{\mathrm{ACV}_{\mathrm{total},t}}
\end{equation*}
\textit{An LTV ratio below $1.0$ indicates collateral value exceeds the exposure.
An LTV ratio at or above $1.0$ triggers action.}
\end{formulabox}

\begin{tabularx}{\linewidth}{lX}
\toprule
\vartableheader
\midrule
$\mathrm{NLV}_t$                   & Net Liquidation Value at time $t$ \\
$\mathrm{MV}_{\mathrm{portfolio},t}$ & Market value of all open positions at time $t$ \\
$L_t$                              & Outstanding loan balance at time $t$ \\
$\mathrm{LTV}_t$                   & Loan-to-Value ratio at time $t$ \\
$\mathrm{ACV}_{\mathrm{total},t}$  & Total Accepted Collateral Value at time $t$
                                     ($= \mathrm{NLV}_t + \text{additional posted collateral}$) \\
\bottomrule
\end{tabularx}

% ─────────────────────────────────────────────────────────
\subsubsection{E.3.3 — Stress Testing LTV}

In addition to the current LTV ratio, a stress LTV is calculated daily by
applying the stressed scenario prices from the initial margin framework to both
the financed positions and the collateral. It identifies positions at material
risk of breaching margin breach limits under adverse market conditions.

\begin{formulabox}
\textbf{Stress LTV}
\begin{equation*}
    \mathrm{LTV}_{\mathrm{stress}} =
        \frac{L_t}{\mathrm{ACV}_{\mathrm{stress}}}
\end{equation*}
\end{formulabox}

\begin{formulabox}
\textbf{Stressed Accepted Collateral Value}
\begin{equation*}
    \mathrm{ACV}_{\mathrm{stress}} =
        \sum_{i} \mathrm{MV}_{\mathrm{stress}}(i)
        \times \bigl(1 - H_{\mathrm{eff}}(i)\bigr)
        + \sum_{j} \mathrm{Cash}(j)
\end{equation*}
\textit{$\mathrm{MV}_{\mathrm{stress}}(i)$ uses scenario prices from the
stressed scenario set of the initial margin framework.}
\end{formulabox}

\begin{tabularx}{\linewidth}{lX}
\toprule
\vartableheader
\midrule
$\mathrm{LTV}_{\mathrm{stress}}$     & Stress Loan-to-Value ratio \\
$\mathrm{ACV}_{\mathrm{stress}}$     & Accepted Collateral Value under stress scenario prices \\
$\mathrm{MV}_{\mathrm{stress}}(i)$  & Market value of collateral instrument $i$
                                       at stressed scenario prices \\
\bottomrule
\end{tabularx}

\newpage

% ─────────────────────────────────────────────────────────
%  APPENDIX — CONSOLIDATED FORMULA REFERENCE
% ─────────────────────────────────────────────────────────
\appendix
\section{Consolidated Formula Reference}
\label{app:formulas}

This appendix lists all quantitative formulas from the methodology in the order
in which they appear, with their canonical names.

\renewcommand{\arraystretch}{1.6}

\begin{longtable}{p{0.05\linewidth} p{0.35\linewidth} p{0.52\linewidth}}
\toprule
\textbf{Ref.} & \textbf{Name} & \textbf{Formula} \\
\midrule
\endfirsthead
\toprule
\textbf{Ref.} & \textbf{Name} & \textbf{Formula} \\
\midrule
\endhead
\midrule
\multicolumn{3}{r}{\small\itshape Continued on next page}\\
\endfoot
\bottomrule
\endlastfoot

%% ── Section D: Margin ──────────────────────────────────

\multicolumn{3}{l}{\textbf{D.1 — Variation Margin}} \\[2pt]

1
    & Variation Margin
    & $\displaystyle \mathrm{VM}_{i,t} = \bigl(P_{i,t} - P_{i,t-1}\bigr) \times N_i$ \\[6pt]

%% ── D.2.2 Scenario Construction ────────────────────────

\multicolumn{3}{l}{\textbf{D.2.2 — Scenario Construction}} \\[2pt]

2
    & Filtered (Volatility-Scaled) Return
    & $\displaystyle r_{\mathrm{filt}}(i,s) = r_{\mathrm{hist}}(i,s)
       \times \dfrac{\hat{\sigma}_{i,0}}{\hat{\sigma}_{i,s}}$ \\[8pt]

3
    & Filtered Historical Scenario Price
    & $\displaystyle P_{\mathrm{filt}}(i,s) = P_{i,0}
       \times \bigl(1 + r_{\mathrm{filt}}(i,s)\bigr)$ \\[6pt]

4
    & Stressed Scenario Price
    & $\displaystyle P_{\mathrm{stress}}(i,s) = P_{i,0}
       \times \bigl(1 + r_{\mathrm{stress}}(i,s)\bigr)$ \\[6pt]

%% ── D.2.3 P&L Distributions ────────────────────────────

\multicolumn{3}{l}{\textbf{D.2.3 — Profit and Loss Distributions}} \\[2pt]

5
    & Instrument P\&L
    & $\displaystyle \mathrm{PL}(i,s) =
       \bigl(P_{i,0} - P_{\mathrm{scenario}}(i,s)\bigr) \times N_i$ \\[6pt]

6
    & Liquidation Group Split P\&L
    & $\displaystyle \mathrm{PL}(l,s) = \sum_{i \,\in\, l} \mathrm{PL}(i,s)$ \\[6pt]

%% ── D.2.4.1 Market Risk Model ──────────────────────────

\multicolumn{3}{l}{\textbf{D.2.4.1 — Market Risk Model}} \\[2pt]

7
    & Empirical Quantile
    & $\displaystyle q(i) = 100\% \times \dfrac{i-0.5}{n}$ \\[8pt]

8
    & Interpolated VaR
    & $\displaystyle \mathrm{VaR}(\alpha) = -\left[
       \dfrac{\mathrm{PL}(q^{+})(q^{+}-(1-\alpha){\cdot}100)
            + \mathrm{PL}(q^{-})((1-\alpha){\cdot}100 - q^{-})}
             {q^{+} - q^{-}}
      \right]$ \\[8pt]

9
    & Robust VaR Scaling Factor
    & $\displaystyle \mathrm{SF} =
       \dfrac{t_{\nu}^{-1}(\alpha_{\mathrm{high}})}
             {t_{\nu}^{-1}(\alpha_{\mathrm{low}})}$ \\[8pt]

10
    & Robust VaR
    & $\displaystyle \mathrm{RVaR}(l,a) =
       \mathrm{VaR}(\alpha_{\mathrm{low}},l,a) \times \mathrm{SF}$ \\[6pt]

11
    & Market Risk Component
    & $\displaystyle \mathrm{MRC}(l) = \max\!\Bigl(
       \overline{\mathrm{RVaR}}_{\mathrm{hist}}(l) + \overline{\mathrm{CB}}(l),\;
       w_s \cdot \overline{\mathrm{RVaR}}_{\mathrm{stress}}(l)
      \Bigr)$ \\[6pt]

%% ── D.2.4.2 Model Risk Buffer ──────────────────────────

\multicolumn{3}{l}{\textbf{D.2.4.2 — Model Risk Buffer (Correlation Break Adjustment)}} \\[2pt]

12
    & Excess Risk per Sub-Window
    & $\displaystyle \mathrm{RVaR}^{+}(l,a,T_j) =
       \max\!\bigl(0,\,
         \mathrm{RVaR}_{\mathrm{sub}}(l,a,T_j) - \mathrm{RVaR}_{\mathrm{hist}}(l,a)
       \bigr)$ \\[6pt]

13
    & Root Mean Square Excess Risk
    & $\displaystyle \mathrm{RMSE}(l,a) =
       \sqrt{\dfrac{1}{n^{+}_a}
             \sum_{j=1}^{n^{+}_a}
             \bigl[\mathrm{RVaR}^{+}(l,a,T_j)\bigr]^{2}}$ \\[8pt]

14
    & Correlation Break Adjustment
    & $\displaystyle \mathrm{CB}(l,a) = \mathrm{clip}\!\Bigl(
       \kappa \cdot \mathrm{RMSE}(l,a),\;
       \Phi \cdot \mathrm{RVaR}_{\mathrm{hist}}(l,a),\;
       \psi \cdot \mathrm{RVaR}_{\mathrm{hist}}(l,a)
      \Bigr)$ \\[6pt]

%% ── D.2.5 Liquidity Risk Component ─────────────────────

\multicolumn{3}{l}{\textbf{D.2.5 — Liquidity Risk Component}} \\[2pt]

15
    & Position-Level Liquidity Charge
    & $\displaystyle \mathrm{LC}(i) =
       \mathrm{RVaR}_{\mathrm{pos}}(i) \times \lambda(i)$ \\[6pt]

16
    & Diversification Factor
    & $\displaystyle \mathrm{DF}(l) =
       \max\!\left(
         \mathrm{DF}_{\mathrm{floor}},\;
         \dfrac{\mathrm{RVaR}_{\mathrm{portfolio}}(l)}
               {\sum_{i \,\in\, l} \mathrm{RVaR}_{\mathrm{pos}}(i)}
       \right)$ \\[8pt]

17
    & Aggregate Liquidity Risk Component
    & $\displaystyle \mathrm{LC}(l) =
       \mathrm{DF}(l) \times \sum_{i \,\in\, l} \mathrm{LC}(i)$ \\[6pt]

%% ── D.2.6 Initial Margin ───────────────────────────────

\multicolumn{3}{l}{\textbf{D.2.6 — Aggregation and Final Initial Margin}} \\[2pt]

18
    & Initial Margin (Split)
    & $\displaystyle \mathrm{IM}(l) = \mathrm{MRC}(l) + \mathrm{LC}(l)$ \\[6pt]

19
    & Total Initial Margin
    & $\displaystyle \mathrm{IM}_{\mathrm{total}} = \sum_{l} \mathrm{IM}(l)$ \\[6pt]

%% ── E.2 Collateral LTV ─────────────────────────────────

\multicolumn{3}{l}{\textbf{E.2.2 — Haircut Calibration}} \\[2pt]

20
    & Collateral LTV
    & $\displaystyle \mathrm{LTV}_{\mathrm{coll}}(i) = 1 - H(i)$ \\[6pt]

21
    & Accepted Collateral Value
    & $\displaystyle \mathrm{ACV}(i) = \mathrm{MV}(i) \times \bigl(1 - H(i)\bigr)$ \\[6pt]

22
    & Total Accepted Collateral Value
    & $\displaystyle \mathrm{ACV}_{\mathrm{total}} =
       \sum_{i} \mathrm{ACV}(i) + \sum_{j} \mathrm{Cash}(j)$ \\[6pt]

%% ── E.2.3 Concentration Add-ons ────────────────────────

\multicolumn{3}{l}{\textbf{E.2.3 — Concentration Add-ons}} \\[2pt]

23
    & Average Pairwise Correlation
    & $\displaystyle \bar{\rho}(g) =
       \dfrac{2}{m(m-1)}
       \sum_{i < j,\; i,j\,\in\, g} \rho(i,j)$ \\[8pt]

24
    & Diversification Score
    & $\displaystyle \mathrm{DS}(g) =
       \sqrt{\dfrac{1 + (m-1)\,\bar{\rho}(g)}{m}}$ \\[8pt]

25
    & Concentration Add-on Haircut
    & $\displaystyle H_{\mathrm{conc}}(i,g) =
       H_{\mathrm{base}}(i)
       \times \dfrac{\mathrm{DS}(g) - \mathrm{DS}_{\mathrm{floor}}}
                    {1 - \mathrm{DS}_{\mathrm{floor}}}$ \\[8pt]

26
    & Effective Haircut
    & $\displaystyle H_{\mathrm{eff}}(i) =
       H_{\mathrm{base}}(i) + H_{\mathrm{conc}}(i,g)$ \\[6pt]

%% ── E.3 LTV ────────────────────────────────────────────

\multicolumn{3}{l}{\textbf{E.3 — Loan-to-Value for Margin Lending}} \\[2pt]

27
    & Net Liquidation Value
    & $\displaystyle \mathrm{NLV}_t =
       \mathrm{MV}_{\mathrm{portfolio},t} - L_t$ \\[6pt]

28
    & LTV Ratio
    & $\displaystyle \mathrm{LTV}_t =
       \dfrac{L_t}{\mathrm{ACV}_{\mathrm{total},t}}$ \\[8pt]

29
    & Stress LTV
    & $\displaystyle \mathrm{LTV}_{\mathrm{stress}} =
       \dfrac{L_t}{\mathrm{ACV}_{\mathrm{stress}}}$ \\[6pt]

30
    & Stressed Accepted Collateral Value
    & $\displaystyle \mathrm{ACV}_{\mathrm{stress}} =
       \sum_{i} \mathrm{MV}_{\mathrm{stress}}(i)
       \times (1 - H_{\mathrm{eff}}(i))
       + \sum_{j} \mathrm{Cash}(j)$ \\[6pt]

\end{longtable}

\renewcommand{\arraystretch}{1.0}

\vfill
\begin{center}
    \rule{0.5\linewidth}{0.4pt}\\[0.5em]
    {\small\itshape Strictly Private and Confidential.
    This document is intended for internal use by Trade Republic's risk management
    function and authorised counterparties. It does not constitute financial or
    legal advice. All margin calculations are subject to applicable rules,
    regulations, and the Trade Republic margin trading eligibility policy.}
\end{center}

% ─────────────────────────────────────────────────────────
\section{Worked Example — Customer Opens a 3$\times$ Leveraged Position}
% ─────────────────────────────────────────────────────────

The following example illustrates the end-to-end margin calculation for a customer
wishing to open a 3$\times$ leveraged position in an equity ETF.

\subsection*{Mock Data}

\begin{tabularx}{\linewidth}{lX}
\toprule
\textbf{Parameter} & \textbf{Value} \\
\midrule
Instrument          & MSCI World ETF (IE00B4L5Y983) \\
Current price $P_0$ & \EUR{100.00} per share \\
Quantity $N$        & 150 shares \\
Position value      & $\mathrm{MV}_0 = 150 \times \EUR{100} = \EUR{15{,}000}$ \\
Target leverage     & $3\times$ \\
Customer equity     & $\EUR{15{,}000} / 3 = \EUR{5{,}000}$ \\
Loan $L_0$          & $\EUR{15{,}000} - \EUR{5{,}000} = \EUR{10{,}000}$ \\
Annualised vol $\sigma$ & $18\%$ \\
Asset-class haircut floor $h_{\text{floor}}$ & $10\%$ (equity ETF) \\
Volatility buffer weight $w_{\text{buf}}$ & $0.50$ \\
Haircut horizon     & 1 trading day \\
\bottomrule
\end{tabularx}

\bigskip

% ── Step 1 ──────────────────────────────────────────────
\subsection*{Step 1 — Collateral Assessment (Phase 1 / Section E.2)}

\textbf{Volatility-implied haircut} over a 1-day horizon at 99\%:
\[
    h_{\text{vol}} = 1 - e^{-z_{0.99}\,\sigma\,\sqrt{1/252}}
                   = 1 - e^{-2.326 \times 0.18 \times 0.0630}
                   \approx 2.6\%
\]
Since $h_{\text{vol}} = 2.6\% < h_{\text{floor}} = 10\%$, the floor binds and the vol
excess is zero:
\[
    h_{\text{base}} = 10\% + 0.5 \times \max(0,\; 2.6\% - 10\%) = \mathbf{10\%}
\]
No concentration add-on ($h_{\text{conc}} = 0$), so $h_{\text{eff}} = 10\%$.

\[
    \mathrm{ACV}_0 = \EUR{15{,}000} \times (1 - 0.10) = \mathbf{\EUR{13{,}500}}
\]

% ── Step 2 ──────────────────────────────────────────────
\subsection*{Step 2 — Initial Margin Calculation (Phase 2 / Section D.2)}

From 250 filtered historical scenarios (volatility-scaled, no Student-t anchor):
\[
    \mathrm{RVaR}^{\text{hist}} \approx 5.9\% \quad
    \text{(99th percentile loss over 5-day liquidation horizon)}
\]
\[
    \mathrm{CB} = 0.5\% \quad \Rightarrow \quad
    \mathrm{MRC} = 5.9\% + 0.5\% = \mathbf{6.4\%}
\]
\[
    \mathrm{LC} = \mathrm{DF} \times \mathrm{RVaR}^{\text{hist}} \times \lambda
               = 0.85 \times 5.9\% \times 1.0 = \mathbf{5.0\%}
\]
\[
    \mathrm{IM}_{\text{rate}} = 6.4\% + 5.0\% = \mathbf{11.4\%}
    \quad \Rightarrow \quad
    \mathrm{IM}_{\text{EUR}} = 11.4\% \times \EUR{15{,}000} = \mathbf{\EUR{1{,}710}}
\]

% ── Step 3 ──────────────────────────────────────────────
\subsection*{Step 3 — Trade Execution Gate (Phase 3)}

\begin{tabularx}{\linewidth}{lllX}
\toprule
\textbf{Check} & \textbf{Value} & \textbf{Threshold} & \textbf{Result} \\
\midrule
Customer equity vs.\ IM & \EUR{5{,}000} $\geq$ \EUR{1{,}710} & IM$_{\text{EUR}}$ & \textbf{Pass} \\
ACV vs.\ IM             & \EUR{13{,}500} $\geq$ \EUR{1{,}710} & IM$_{\text{EUR}}$ & \textbf{Pass} \\
Opening LTV             & $\EUR{10{,}000} / \EUR{13{,}500} = 74.1\%$ & $75\%$ & \textbf{Pass} \\
\bottomrule
\end{tabularx}

\medskip
All checks pass. \textbf{Position is opened.}

\[
    \mathrm{LTV}_0 = \frac{L_0}{\mathrm{ACV}_0}
                   = \frac{\EUR{10{,}000}}{\EUR{13{,}500}} = 74.1\%
\]

\[
    \text{Max leverage at LTV}=75\%:\quad
    L^* = \frac{1}{0.25 + 0.75 \times 0.10} = \frac{1}{0.325} = 3.08\times
    \quad \checkmark
\]

% ── Step 4a ─────────────────────────────────────────────
\subsection*{Step 4a — Day 1: Adverse Price Move (Phase 4 / Section D.1.3)}

Price falls from \EUR{100} to \EUR{97} ($-3\%$):

\[
    \mathrm{VM}_1 = (97 - 100) \times 150 = -\EUR{450}
    \quad \text{(debited from customer account)}
\]
\[
    \mathrm{MV}_1 = 97 \times 150 = \EUR{14{,}550}
    \qquad
    \mathrm{ACV}_1 = \EUR{14{,}550} \times 0.90 = \EUR{13{,}095}
\]
\[
    \mathrm{LTV}_1 = \frac{\EUR{10{,}000}}{\EUR{13{,}095}} = \mathbf{76.4\%}
    > 75\% \quad \Rightarrow \quad \textbf{Monitoring alert triggered}
\]

% ── Step 5a ─────────────────────────────────────────────
\subsection*{Step 5a — Day 1: LTV Limit Check (Phase 5 / Section G.3.4)}

\begin{tabularx}{\linewidth}{llll}
\toprule
\textbf{Threshold} & \textbf{Limit} & \textbf{LTV$_1$} & \textbf{Action} \\
\midrule
Monitoring  & $75\%$ & $76.4\%$ & \textbf{Flag — enhanced surveillance} \\
Margin breach & $80\%$ & $76.4\%$ & No action \\
Liquidation & $90\%$ & $76.4\%$ & No action \\
\bottomrule
\end{tabularx}

% ── Step 4b ─────────────────────────────────────────────
\subsection*{Step 4b — Day 2: Price Recovery (Phase 4)}

Price recovers from \EUR{97} to \EUR{99} ($+2.1\%$):

\[
    \mathrm{VM}_2 = (99 - 97) \times 150 = +\EUR{300}
    \quad \text{(credited to customer account)}
\]
\[
    \mathrm{MV}_2 = 99 \times 150 = \EUR{14{,}850}
    \qquad
    \mathrm{ACV}_2 = \EUR{14{,}850} \times 0.90 = \EUR{13{,}365}
\]
\[
    \mathrm{LTV}_2 = \frac{\EUR{10{,}000}}{\EUR{13{,}365}} = \mathbf{74.8\%}
    < 75\% \quad \Rightarrow \quad \textbf{Monitoring flag cleared — position healthy}
\]

% ── Summary ─────────────────────────────────────────────
\subsection*{Summary}

\begin{tabularx}{\linewidth}{lllll}
\toprule
\textbf{Day} & \textbf{Price} & \textbf{VM} & \textbf{ACV} & \textbf{LTV} \\
\midrule
$t=0$ (open) & \EUR{100.00} & --- & \EUR{13{,}500} & $74.1\%$ \quad OK \\
$t=1$        & \EUR{97.00}  & $-$\EUR{450} & \EUR{13{,}095} & $76.4\%$ \quad \textit{Monitoring} \\
$t=2$        & \EUR{99.00}  & $+$\EUR{300} & \EUR{13{,}365} & $74.8\%$ \quad OK \\
\bottomrule
\end{tabularx}

\bigskip
The example demonstrates that a $3\times$ leveraged position in this instrument opens
comfortably within the 75\% LTV limit and that daily Variation Margin adjustments
provide a dynamic early-warning mechanism. A single adverse 3\% daily move is
sufficient to trigger the monitoring threshold, highlighting the sensitivity of
leveraged positions to price changes and the importance of the ongoing maintenance
margin framework.

% ─────────────────────────────────────────────────────────
\subsection*{Extended Example — Monitoring Framework Across Liquidation Horizon}
% ─────────────────────────────────────────────────────────

The following extends the worked example above to illustrate the two-phase monitoring
framework (Section~D.1.5), including the loan evolution through Variation Margin,
the overnight haircut check, and the transition to the seasoned phase at $t = 5$.

\paragraph{Setup.}
\begin{tabularx}{\linewidth}{lX}
\toprule
\textbf{Parameter} & \textbf{Value} \\
\midrule
Opening price $P_0$              & \EUR{100.00} \\
Quantity $N$                     & 150 shares \\
Position value $\mathrm{MV}_0$   & \EUR{15{,}000} \\
Opening loan $L_0$               & \EUR{10{,}000} \\
Intraday haircut $h_{\text{eff}}$     & $10\%$ \\
Overnight haircut $h_{\text{overnight}}$ & $14\%$ \\
Monitoring threshold $\overline{\mathrm{LTV}}$ & $75\%$ \\
Liquidation threshold $\mathrm{LTV}_{\text{liq}}$ & $90\%$ \\
Liquidation horizon $T_{\text{liq}}$ & 4 trading days \\
\bottomrule
\end{tabularx}

\bigskip

\paragraph{Step 1 — Collateral Assessment.}
\[
    h_{\text{vol}} = 1 - e^{-2.326 \times 0.18 \times \sqrt{1/252}} \approx 2.6\%
    \;<\; h_{\text{floor}} = 10\% \quad\Rightarrow\quad h_{\text{eff}} = 10\%
\]
\[
    \mathrm{ACV}_0 = \EUR{15{,}000} \times (1 - 0.10) = \EUR{13{,}500}
\]

\paragraph{Step 2 — Initial Margin.}
\[
    \mathrm{RVaR}^{\text{hist}} \approx 5.9\%,\quad
    \mathrm{CB} = 0.5\%,\quad
    \mathrm{MRC} = 6.4\%,\quad
    \mathrm{LC} = 5.0\%
\]
\[
    \mathrm{IM}_{\text{rate}} = 11.4\%
    \quad\Rightarrow\quad
    \mathrm{IM}_{\text{EUR}} = \EUR{1{,}710}
\]

\paragraph{Step 3 — Execution Gate.}
\[
    \mathrm{LTV}_0 = \frac{\EUR{10{,}000}}{\EUR{13{,}500}} = 74.1\%
    \;<\; 75\% \quad\checkmark \quad \textbf{Position opened.}
\]

\paragraph{Step 4 — Day 1: Price $-3\%$ to \EUR{97} \textnormal{(Phase 1 — within liquidation horizon)}.}
\[
    \mathrm{VM}_1 = (97 - 100) \times 150 = -\EUR{450}
    \qquad
    L_1 = \EUR{10{,}000} - (-\EUR{450}) = \EUR{10{,}450}
\]
\[
    \mathrm{ACV}_1 = \EUR{14{,}550} \times 0.90 = \EUR{13{,}095}
    \qquad
    \mathrm{LTV}_1 = \frac{\EUR{10{,}450}}{\EUR{13{,}095}} = 79.8\%
    \;>\; 75\%
\]
\textbf{Monitoring alert triggered (IM\,+\,VM breach, Phase~1).}

\paragraph{Step 5 — Day 1 End-of-Day: Overnight Check.}
\[
    \mathrm{ACV}_1^{\text{overnight}} = \EUR{14{,}550} \times (1 - 0.14) = \EUR{12{,}513}
    \qquad
    \mathrm{LTV}_1^{\text{overnight}} = \frac{\EUR{10{,}450}}{\EUR{12{,}513}} = 83.5\%
    \;>\; 75\%
\]
\textbf{Overnight margin review triggered.} Customer must post additional collateral
or reduce the loan before next open. $83.5\% < 90\%$: no liquidation.

\paragraph{Step 6 — Day 2: Price $+2\%$ to \EUR{99}.}
\[
    \mathrm{VM}_2 = (99 - 97) \times 150 = +\EUR{300}
    \qquad
    L_2 = \EUR{10{,}450} - \EUR{300} = \EUR{10{,}150}
\]
\[
    \mathrm{ACV}_2 = \EUR{14{,}850} \times 0.90 = \EUR{13{,}365}
    \qquad
    \mathrm{LTV}_2 = \frac{\EUR{10{,}150}}{\EUR{13{,}365}} = 75.9\%
    \;>\; 75\%
\]
Monitoring flag remains active; no overnight breach ($\mathrm{LTV}^{\text{overnight}}_2 < 75\%$ if
overnight haircut does not push ACV below threshold).

\paragraph{Step 7 — Day 5: Position Seasoned ($t > T_{\text{liq}}$), Price stable at \EUR{99}.}
Phase~2 now applies. Initial Margin is no longer binding; Maintenance Margin governs.
\[
    \mathrm{LTV}_5 = \frac{\EUR{10{,}150}}{\EUR{13{,}365}} = 75.9\%
    \;>\; 75\%
    \quad\Rightarrow\quad \textbf{MM monitoring alert (Phase~2).}
\]
Customer must top up collateral or reduce the loan to clear the flag.

\bigskip

\paragraph{Summary.}

\begin{tabularx}{\linewidth}{llllll}
\toprule
\textbf{Day} & \textbf{Price} & \textbf{VM} & \textbf{Loan} & \textbf{ACV} & \textbf{LTV / Status} \\
\midrule
$t=0$ (open) & \EUR{100} & ---         & \EUR{10{,}000} & \EUR{13{,}500} & $74.1\%$ \quad OK \\
$t=1$        & \EUR{97}  & $-$\EUR{450} & \EUR{10{,}450} & \EUR{13{,}095} & $79.8\%$ \quad \textit{Monitoring (IM+VM)} \\
$t=1$ EOD    & \EUR{97}  & ---          & \EUR{10{,}450} & \EUR{12{,}513}$^\dagger$ & $83.5\%$ \quad \textit{Overnight review} \\
$t=2$        & \EUR{99}  & $+$\EUR{300} & \EUR{10{,}150} & \EUR{13{,}365} & $75.9\%$ \quad \textit{Monitoring} \\
$t=5+$       & \EUR{99}  & ---          & \EUR{10{,}150} & \EUR{13{,}365} & $75.9\%$ \quad \textit{MM monitoring (Phase~2)} \\
\bottomrule
\end{tabularx}

\smallskip
{\small $^\dagger$Overnight ACV uses $h_{\text{overnight}} = 14\%$.}

% ─────────────────────────────────────────────────────────
\newpage
\section*{Appendix: Comparison with Interactive Brokers Margin Framework}
% ─────────────────────────────────────────────────────────

This appendix applies the Interactive Brokers (IB) margin framework to the identical
worked example used in the preceding sections — 150~shares at \EUR{100} per share,
yielding a \EUR{15{,}000} position with \EUR{5{,}000} customer equity and a
\EUR{10{,}000} loan (3$\times$ leverage) — to facilitate a direct methodological comparison.

% ─────────────────────────────────────────────────────────
\subsection*{IB Approach A: Regulation T (Reg~T)}
% ─────────────────────────────────────────────────────────

Reg~T is a flat-rate rule issued by the US Federal Reserve (12~CFR~Part~220) and
enforced by FINRA Rule~4210. It defines a minimum initial margin of 50\% of the
purchase price for long equity positions, irrespective of the actual risk profile
of the instrument.

\paragraph{Step 1 — Compute the Reg~T initial margin requirement.}
\[
    \mathrm{IM}_{\mathrm{RegT}} = 50\% \times \EUR{15{,}000} = \EUR{7{,}500}
\]

\paragraph{Step 2 — Assess whether the position can be opened.}
Customer equity is \EUR{5{,}000}, which is less than the required \EUR{7{,}500}.
\textbf{The \EUR{15{,}000} position cannot be opened under Reg~T.}
The same \EUR{5{,}000} of equity permits a maximum position of:
\[
    \mathrm{MV}_{\max} = \frac{\EUR{5{,}000}}{50\%} = \EUR{10{,}000}
    \quad\Longrightarrow\quad 100 \text{ shares at } \EUR{100} \quad (2{\times}\text{ leverage only})
\]

\paragraph{Step 3 — Adjusted position (maximum under Reg~T).}

\begin{tabularx}{\linewidth}{lX}
\toprule
\textbf{Item} & \textbf{Value} \\
\midrule
Shares & 100 \\
Market value & \EUR{10{,}000} \\
Loan & \EUR{5{,}000} \\
Customer equity & \EUR{5{,}000} \\
Leverage & $2{\times}$ \\
Reg~T initial margin rate & $50\%$ \\
\bottomrule
\end{tabularx}

\paragraph{Step 4 — Maintenance margin requirement (ongoing).}
FINRA Rule~4210 sets a minimum maintenance margin of 25\% of current market value.
IB applies this floor without further add-ons for standard large-cap equities:
\[
    \mathrm{MMR} = 25\% \times \EUR{10{,}000} = \EUR{2{,}500}
\]
\[
    \text{Excess Liquidity} = \text{Equity} - \mathrm{MMR} = \EUR{5{,}000} - \EUR{2{,}500} = \EUR{2{,}500}
\]
IB begins forced liquidation when Excess Liquidity drops below zero, without issuing a
margin call. There is no graduated warning framework.

\paragraph{Step 5 — Compute the auto-liquidation price.}
Let $m = 0.25$ (maintenance rate), $L = \EUR{5{,}000}$, $N = 100$ shares:
\[
    P^{*} = \frac{L}{N \times (1 - m)} = \frac{\EUR{5{,}000}}{100 \times 0.75} = \EUR{66.67}
\]
The stock must fall $-33.3\%$ from \EUR{100} before IB triggers auto-liquidation.

\paragraph{Step 6 — IB account panel (at position open).}

\begin{tabularx}{\linewidth}{lll}
\toprule
\textbf{IB Field} & \textbf{Formula} & \textbf{Value} \\
\midrule
Net Liquidation Value  & Equity                   & \EUR{5{,}000} \\
Initial Margin Req.    & $50\% \times \mathrm{MV}$ & \EUR{5{,}000} \\
Maintenance Margin Req.& $25\% \times \mathrm{MV}$ & \EUR{2{,}500} \\
Available Funds        & Equity $-$ IMR            & \EUR{0} \quad (at limit) \\
Excess Liquidity       & Equity $-$ MMR            & \EUR{2{,}500} \\
Cushion                & Excess Liq.\ / NLV        & $50\%$ \\
\bottomrule
\end{tabularx}

% ─────────────────────────────────────────────────────────
\subsection*{IB Approach B: Portfolio Margin (PM)}
% ─────────────────────────────────────────────────────────

Portfolio Margin is available to IB accounts with net liquidation value above
\$110{,}000 (FINRA Rule~4210(g)). It replaces the flat Reg~T rule with a scenario-based
engine derived from the OCC's Theoretical Intermarket Margin System (TIMS).

\paragraph{Step 1 — Run the OCC scenario grid.}
IB applies 31 price-shock scenarios ranging from $-15\%$ to $+15\%$ in 1\% steps,
combined with implied-volatility shocks of $\pm 75\%$ at each price point.
For a plain long equity position the dominant scenario is the $-15\%$ price shock:
\[
    \text{Scenario loss} = 15\% \times \EUR{15{,}000} = \EUR{2{,}250}
\]

\paragraph{Step 2 — Apply the single-stock PM floor.}
FINRA Rule~4210(g) mandates a minimum requirement of 15\% of market value for
individual equity positions:
\[
    \mathrm{PM}_{\mathrm{floor}} = 15\% \times \EUR{15{,}000} = \EUR{2{,}250}
\]
\[
    \mathrm{PM\_req} = \max(\text{scenario loss},\; \mathrm{PM}_{\mathrm{floor}}) = \EUR{2{,}250}
\]

\paragraph{Step 3 — Assess whether the position can be opened.}
Customer equity \EUR{5{,}000} $>$ \EUR{2{,}250} required.
\textbf{The \EUR{15{,}000} position can be opened under PM (3$\times$ leverage is feasible).}

\paragraph{Step 4 — Maximum leverage under PM.}
\[
    \text{Max leverage} = \frac{1}{15\%} \approx 6.7{\times}
    \quad\Longrightarrow\quad
    \mathrm{MV}_{\max} = \frac{\EUR{5{,}000}}{15\%} = \EUR{33{,}333}
\]

\paragraph{Step 5 — PM maintenance and liquidation trigger.}
The scenario grid is re-evaluated continuously throughout the trading day.
For single equities, the maintenance requirement equals the PM requirement (15\% floor).
Forced liquidation is triggered when Excess Liquidity (NLV $-$ PM\_req) drops below zero.

% ─────────────────────────────────────────────────────────
\subsection*{Side-by-Side Comparison}
% ─────────────────────────────────────────────────────────

The table below compares all three frameworks on the \EUR{15{,}000} worked example.
Where IB~Reg~T requires an adjusted 2$\times$ position, figures are shown on a
comparable per-EUR basis.

\begin{tabularx}{\linewidth}{lXXX}
\toprule
\textbf{Dimension} & \textbf{IB Reg~T} & \textbf{IB Portfolio Margin} & \textbf{Trade Republic} \\
\midrule
Methodology
  & Flat regulatory rule (50\%)
  & OCC scenario grid $\pm 15\%$, 31 points
  & Filtered historical simulation + EWMA volatility scaling \\
Initial margin (€15,000 position)
  & \EUR{7{,}500} (50\%)
  & \EUR{2{,}250} (15\%)
  & \EUR{1{,}710} (11.4\%) \\
3$\times$ leverage possible?
  & \textbf{No} — needs \EUR{7{,}500} equity
  & \textbf{Yes}
  & \textbf{Yes} \\
Max leverage on \EUR{5{,}000} equity
  & $2{\times}$
  & $6.7{\times}$
  & ${\sim}8{\times}$ (LTV-capped at 75\%) \\
IM rate driver
  & Regulatory flat rate
  & Worst-case $\pm 15\%$ scenario + 15\% floor
  & 99\% RVaR over 5-trading-day horizon \\
Maintenance trigger
  & Excess Liquidity $< 0$ (binary)
  & Excess Liquidity $< 0$ (binary)
  & LTV $> 90\%$ (liquidation); LTV $> 75\%$ (monitoring) \\
Graduated warning framework
  & None
  & None
  & Yes — three thresholds (75\% / 80\% / 90\%) \\
Auto-liquidation price (3$\times$, \EUR{100})
  & N/A (not permitted)
  & $P^{*} \approx \EUR{83}$ at 15\% MM
  & LTV $= 90\%$ $\Rightarrow$ $P^{*} \approx \EUR{74}$ \\
Concentration add-on
  & ADV-based escalation (not published)
  & OCC concentration rules
  & $h_{\mathrm{conc}} + 5\%$ when weight $> 20\%$ \\
Liquidation horizon
  & No explicit horizon
  & 1-day (intraday scenarios)
  & 5 trading days \\
Intraday vs.\ overnight distinction
  & Yes — overnight reverts to 50\% Reg~T
  & Yes — overnight PM re-evaluated at 15:50~ET
  & Yes — overnight haircut add-on applied at EOD \\
Margin call issued
  & No — auto-liquidates
  & No — auto-liquidates
  & No — auto-liquidates \\
\bottomrule
\end{tabularx}

\bigskip

\subsection*{Key Observations}

\begin{enumerate}[leftmargin=2em]

\item \textbf{Reg~T is purely rule-based.} The 50\% flat rate bears no relation to the
      actual volatility or liquidity of the instrument. A low-vol blue-chip and a
      speculative small-cap receive identical treatment. This blocks 3$\times$ leverage
      entirely, regardless of the risk profile.

\item \textbf{IB Portfolio Margin is methodologically closest to the Trade Republic
      framework.} Both apply scenario-based stress tests. IB uses the OCC's fixed
      $\pm 15\%$ grid; Trade Republic uses filtered historical simulation calibrated
      to the realised tail of each ISIN over the trailing 750 trading days.

\item \textbf{The Trade Republic IM is lower than the IB PM floor} (11.4\% vs.\ 15\%)
      because the historical approach captures the actual volatility of the stock rather
      than applying a fixed 15\% shock. For a low-volatility instrument the TR rate will
      be materially lower; for a high-volatility instrument it will exceed the IB PM
      floor.

\item \textbf{The liquidation frameworks differ in gradualism.} IB operates a binary
      trigger (Excess Liquidity $< 0$ $\Rightarrow$ liquidate). Trade Republic operates
      a three-level graduated framework (monitoring at 75\%, breach at 80\%, liquidation
      at 90\%), giving the customer explicit visibility and response time before forced
      liquidation occurs.

\item \textbf{Both IB and Trade Republic auto-liquidate} without issuing a traditional
      margin call. The difference is that Trade Republic's monitoring thresholds
      function as early warnings that a formal margin call framework would otherwise provide.

\end{enumerate}

\end{document}
